% READ2-RECOMP5
\section*{READ2-RECOMP5: rereading collapses the $Q_4$ edge-cover state barrier}

\textbf{STATUS: PROVED general read-$\Delta$ upper bound; PROVED/EXACT $Q_4$ linear-width; COMPUTATIONAL $(1,+17)$ calibration. P versus NP remains OPEN.}

\paragraph{Occurrence lemma.}
Let
\[
f(x)=\bigwedge_j C_j(x)
\]
be a CNF in which each input variable occurs in at most $\Delta$ clauses.  Evaluating the
clauses sequentially gives a deterministic syntactic read-$\Delta$ branching program with at
most
\[
L=\sum_j |C_j|
\]
query nodes: query each clause until one literal is satisfied and reject only if a clause is
exhausted.  Thus low input occurrence is an immediate upper bound against bounded-reread
lower-bound attempts.

\paragraph{$Q_4$ representation.}
The 32 genuine 3-clauses on four Boolean variables correspond to the 32 edges of $Q_4$:
each clause is falsified by exactly the two assignments differing in the omitted variable.
The full clause-presence formula is UNSAT exactly when the selected $Q_4$ edges cover all 16
vertices:
\[
EC_{Q_4}(x)=\bigwedge_{v\in V(Q_4)}\bigvee_{e\ni v}x_e.
\]
Every edge variable appears in exactly two vertex clauses.

\paragraph{Exact read-once width.}
The hypercube pathwidth formula gives
\[
pw(Q_4)=1+1+2+3=7.
\]
An exact closed-set dynamic program finds a linear-width-7 edge layout and rejects width 6,
so
\[
\boxed{lw(Q_4)=7}.
\]
One certified frontier profile is
\[
2,3,4,4,5,5,6,6,6,6,7,7,7,7,7,6,
7,7,7,7,7,7,7,6,6,6,5,4,4,3,2,0.
\]
For edge cover, an ordered read-once state only needs the covered/uncovered status of active
frontier vertices.  The reduced OBDD under this layout has 1773 query nodes and peak width
\[
\boxed{128=2^7}.
\]

\paragraph{Read-two collapse.}
Check the 16 vertex clauses sequentially.  Each block queries the four incident edge bits
until it finds a 1.  The total number of query nodes is
\[
\sum_v \deg(v)=16\cdot4=64=2|E(Q_4)|.
\]
Each edge is queried only in the blocks of its two endpoints, so this is a syntactic read-2
branching program:
\[
\boxed{\mathrm{size}_{\mathrm{read2}}(EC_{Q_4})\le64}.
\]
Relative to the certified width-7 OBDD this is a $96.39\%$ node reduction.

More generally, for every graph $G$,
\[
EC_G=\bigwedge_v\bigvee_{e\ni v}x_e
\]
has a syntactic read-2 branching program with exactly $2|E(G)|$ query nodes.  Therefore
edge-cover functions cannot support superlinear read-2 lower bounds.

\paragraph{Exact-weight-eight promise.}
For an eight-edge subset of $Q_4$,
\[
\text{edge cover}
\iff
\text{perfect matching}
\iff
D(4,8)\text{ is UNSAT}.
\]
Hence the same 64-node deterministic read-2 program, with its terminal labels exchanged,
computes SAT correctly on the exact-weight-eight promise.

\paragraph{$(1,+k)$ calibration.}
A concrete choice of the four local incident-edge orders was verified by exact MILP over all
computation stopping cases.  In this 64-node architecture at most 17 distinct variables are
queried twice on any computation path.  The maxima when the first uncovered vertex occurs at
positions $0,\ldots,15$, followed by the all-covered accepting case, are
\[
0,1,2,2,3,4,6,8,8,9,11,11,13,15,15,17,17.
\]
Thus this explicit program is also a
\[
\boxed{(1,+17)\text{-BP of size }64}.
\]
No optimality claim is made for 17 or for 64.

\paragraph{Interpretation.}
The $Q_4$ laboratory now exhibits an explicit memory--recomputation trade.  Read-once
evaluation retains enough frontier state to avoid asking an edge twice; rereading allows the
machine to discard that information and request the edge again at its other endpoint.  Thus
RECOMP is not a wrapper nuisance here: it genuinely destroys the proposed representation-size
obstruction.

The restricted branching-program literature is consistent with this.  Strong lower bounds are
known for read-once and for specially restricted $(1,+k)$ or bounded-repetition models, but
extending read-once measures to general read-$k$ programs is a recognized difficulty.

\paragraph{Decision.}
Retire $Q_4$ edge cover and the $D(4,8)$ promise as candidates for a lower bound intended to
survive read-two recomputation.  Keep them instead as a finite calibration laboratory for the
Pareto curve
\[
(\text{reread variables},\ \text{retained state},\ \text{DAG size},\ \text{path length}).
\]

\paragraph{Next target.}
Construct that Pareto curve explicitly, or move to a function whose useful representations do
not have low input occurrence and test the Savicky--Zak / Li--McKenzie read-$k$ measures
there.

\begin{thebibliography}{9}
\bibitem{SavickyZak00}
P.~Savick\'y and S.~\v{Z}\'ak,
\emph{A Read-Once Lower Bound and a $(1,+k)$-Hierarchy for Branching Programs},
Theoretical Computer Science 238:347--362, 2000.

\bibitem{ReadSomeTwice97}
\emph{A Lower Bound on Branching Programs Reading Some Bits Twice},
Theoretical Computer Science 172:293--301, 1997.

\bibitem{LiMcKenzie24}
Y.~Li and P.~McKenzie,
\emph{Perspective on Complexity Measures Targeting Read-Once Branching Programs},
Information and Computation 301:105230, 2024.

\bibitem{Thilikos00}
D.~M.~Thilikos et al.,
\emph{Algorithms and Obstructions for Linear-Width and Related Search Parameters},
Discrete Applied Mathematics 105:239--271, 2000.

\bibitem{ChandranKavitha06}
L.~S.~Chandran and T.~Kavitha,
\emph{The Treewidth and Pathwidth of Hypercubes},
Discrete Mathematics 306:359--365, 2006.

\bibitem{KobayashiNakahata20}
Y.~Kobayashi and Y.~Nakahata,
\emph{A Note on Exponential-Time Algorithms for Linearwidth},
arXiv:2010.02388, 2020.
\end{thebibliography}

% END_OF_NOTE::READ2-RECOMP5
